
\documentclass{article}
\usepackage{colortbl}
\usepackage{makecell}
\usepackage{multirow}
\usepackage{supertabular}

\begin{document}

\newcounter{utterance}

\twocolumn

{ \footnotesize  \setcounter{utterance}{1}
\setlength{\tabcolsep}{0pt}
\begin{supertabular}{c@{$\;$}|p{.15\linewidth}@{}p{.15\linewidth}p{.15\linewidth}p{.15\linewidth}p{.15\linewidth}p{.15\linewidth}}

    \# & $\;$A & \multicolumn{4}{c}{Game Master} & $\;\:$B\\
    \hline 

    \theutterance \stepcounter{utterance}  

    & & \multicolumn{4}{p{0.6\linewidth}}{\cellcolor[rgb]{0.9,0.9,0.9}{%
	\makecell[{{p{\linewidth}}}]{% 
	  \tt {\tiny [A$\langle$GM]}  
	 State giocando a un gioco di negoziazione cooperativa in cui dovete accordarvi su come dividere una serie di oggetti tra voi e un altro giocatore.\\ \tt \\ \tt Regole:\\ \tt (a) A Lei e all'altro giocatore viene dato un insieme di oggetti. Ognuno di voi riceve anche una funzione di valore segreta, che rappresenta il valore di ciascun tipo di oggetto.\\ \tt (b) Si scambiano messaggi con l'altro giocatore per concordare chi si aggiudica gli oggetti. Potete inviare un massimo di 5 messaggi ciascuno, oppure terminare in anticipo facendo una proposta segreta in qualsiasi momento.\\ \tt (c) A ciascuno di voi viene chiesto di inviare una proposta segreta indicando gli oggetti che desiderate, formattata tra parentesi quadre come segue: "[Proposta: <numero> <nome oggetto>, <numero> <nome oggetto>, <...>]".\\ \tt (d) Se le vostre proposte sono complementari, cioè ci sono abbastanza oggetti per soddisfare entrambe le proposte, a ciascun giocatore viene assegnato un punteggio basato sulla somma dei valori degli oggetti ricevuti. In caso contrario, entrambi ricevono zero punti.\\ \tt (e) L'obiettivo è massimizzare la somma del proprio punteggio e di quello dell'altro giocatore.\\ \tt \\ \tt Cominciamo.\\ \tt \\ \tt L'insieme degli oggetti disponibili è:\\ \tt \\ \tt 1 anello, 3 frutti, 1 gesso\\ \tt \\ \tt La funzione valore segreta è:\\ \tt \\ \tt anello: 6, frutta: 1, gesso: 1\\ \tt \\ \tt IMPORTANTE: i vostri messaggi, a meno che non si tratti di una proposta segreta, vengono trasmessi direttamente all'altro giocatore, quindi non includete alcuna risposta alle regole o testo di annuncio del vostro messaggio. Per fare una proposta segreta, utilizzate il formato indicato. Non utilizzare delle parentesi quadre quando si comunica all'altro giocatore, altrimenti verrà interpretata come una proposta segreta.\\ \tt  
	  } 
	   } 
	   } 
	 & \\ 
 

    \theutterance \stepcounter{utterance}  

    & \multicolumn{4}{p{0.6\linewidth}}{\cellcolor[rgb]{0.8,1,0.9}{%
	\makecell[{{p{\linewidth}}}]{% 
	  \tt {\tiny [A$\rangle$GM]}  
	 Ciao! Abbiamo a disposizione 1 anello, 3 frutti e 1 gesso da dividere tra noi. Qual è la tua preferenza tra questi oggetti? Io sono particolarmente interessato all'anello, ma sono aperto a discutere diverse opzioni di divisione che possano funzionare bene per entrambi. 
	  } 
	   } 
	   } 
	 & & \\ 
 

    \theutterance \stepcounter{utterance}  

    & & & \multicolumn{2}{p{0.3\linewidth}}{\cellcolor[rgb]{0.95,0.95,0.95}{%
	\makecell[{{p{\linewidth}}}]{% 
	  \tt {\tiny [GM$|$GM]}  
	 continue 
	  } 
	   } 
	   } 
	 & & \\ 
 

    \theutterance \stepcounter{utterance}  

    & & & \multicolumn{2}{p{0.3\linewidth}}{\cellcolor[rgb]{0.95,0.95,0.95}{%
	\makecell[{{p{\linewidth}}}]{% 
	  \tt {\tiny [GM$|$GM]}  
	 Ciao! Abbiamo a disposizione 1 anello, 3 frutti e 1 gesso da dividere tra noi. Qual è la tua preferenza tra questi oggetti? Io sono particolarmente interessato all'anello, ma sono aperto a discutere diverse opzioni di divisione che possano funzionare bene per entrambi. 
	  } 
	   } 
	   } 
	 & & \\ 
 

    \theutterance \stepcounter{utterance}  

    & & \multicolumn{4}{p{0.6\linewidth}}{\cellcolor[rgb]{0.9,0.9,0.9}{%
	\makecell[{{p{\linewidth}}}]{% 
	  \tt {\tiny [GM$\rangle$B]}  
	 State giocando a un gioco di negoziazione cooperativa in cui dovete accordarvi su come dividere una serie di oggetti tra voi e un altro giocatore.\\ \tt \\ \tt Regole:\\ \tt (a) A Lei e all'altro giocatore viene dato un insieme di oggetti. Ognuno di voi riceve anche una funzione di valore segreta, che rappresenta il valore di ciascun tipo di oggetto.\\ \tt (b) Si scambiano messaggi con l'altro giocatore per concordare chi si aggiudica gli oggetti. Potete inviare un massimo di 5 messaggi ciascuno, oppure terminare in anticipo facendo una proposta segreta in qualsiasi momento.\\ \tt (c) A ciascuno di voi viene chiesto di inviare una proposta segreta indicando gli oggetti che desiderate, formattata tra parentesi quadre come segue: "[Proposta: <numero> <nome oggetto>, <numero> <nome oggetto>, <...>]".\\ \tt (d) Se le vostre proposte sono complementari, cioè ci sono abbastanza oggetti per soddisfare entrambe le proposte, a ciascun giocatore viene assegnato un punteggio basato sulla somma dei valori degli oggetti ricevuti. In caso contrario, entrambi ricevono zero punti.\\ \tt (e) L'obiettivo è massimizzare la somma del proprio punteggio e di quello dell'altro giocatore.\\ \tt \\ \tt Cominciamo.\\ \tt \\ \tt L'insieme degli oggetti disponibili è:\\ \tt \\ \tt 1 anello, 3 frutti, 1 gesso\\ \tt \\ \tt La funzione valore segreta è:\\ \tt \\ \tt anello: 6, frutta: 1, gesso: 1\\ \tt \\ \tt IMPORTANTE: i vostri messaggi, a meno che non si tratti di una proposta segreta, vengono trasmessi direttamente all'altro giocatore, quindi non includete alcuna risposta alle regole o testo di annuncio del vostro messaggio. Per fare una proposta segreta, utilizzate il formato indicato. Non utilizzare delle parentesi quadre quando si comunica all'altro giocatore, altrimenti verrà interpretata come una proposta segreta.\\ \tt \\ \tt \\ \tt Ciao! Abbiamo a disposizione 1 anello, 3 frutti e 1 gesso da dividere tra noi. Qual è la tua preferenza tra questi oggetti? Io sono particolarmente interessato all'anello, ma sono aperto a discutere diverse opzioni di divisione che possano funzionare bene per entrambi. 
	  } 
	   } 
	   } 
	 & \\ 
 

    \theutterance \stepcounter{utterance}  

    & & & \multicolumn{4}{p{0.6\linewidth}}{\cellcolor[rgb]{1,0.85,0.72}{%
	\makecell[{{p{\linewidth}}}]{% 
	  \tt {\tiny [GM$\langle$B]}  
	 Ciao! Grazie per aver iniziato la discussione. Anch'io sono interessato all'anello, ma capisco che anche tu lo sia. Per quanto riguarda gli altri oggetti, non ho preferenze particolarmente forti tra frutti e gesso. \\ \tt \\ \tt Dato che entrambi vogliamo l'anello, forse potremmo pensare a una divisione dove chi prende l'anello lascia più degli altri oggetti all'altro? Cosa ne pensi se chi prende l'anello si accontenta di meno frutti, e l'altro prende la maggior parte dei frutti più il gesso?\\ \tt \\ \tt Quali sono le tue preferenze relative tra frutti e gesso? 
	  } 
	   } 
	   } 
	  \\ 
 

    \theutterance \stepcounter{utterance}  

    & & & \multicolumn{2}{p{0.3\linewidth}}{\cellcolor[rgb]{0.95,0.95,0.95}{%
	\makecell[{{p{\linewidth}}}]{% 
	  \tt {\tiny [GM$|$GM]}  
	 continue 
	  } 
	   } 
	   } 
	 & & \\ 
 

    \theutterance \stepcounter{utterance}  

    & & & \multicolumn{2}{p{0.3\linewidth}}{\cellcolor[rgb]{0.95,0.95,0.95}{%
	\makecell[{{p{\linewidth}}}]{% 
	  \tt {\tiny [GM$|$GM]}  
	 Ciao! Grazie per aver iniziato la discussione. Anch'io sono interessato all'anello, ma capisco che anche tu lo sia. Per quanto riguarda gli altri oggetti, non ho preferenze particolarmente forti tra frutti e gesso. \\ \tt \\ \tt Dato che entrambi vogliamo l'anello, forse potremmo pensare a una divisione dove chi prende l'anello lascia più degli altri oggetti all'altro? Cosa ne pensi se chi prende l'anello si accontenta di meno frutti, e l'altro prende la maggior parte dei frutti più il gesso?\\ \tt \\ \tt Quali sono le tue preferenze relative tra frutti e gesso? 
	  } 
	   } 
	   } 
	 & & \\ 
 

    \theutterance \stepcounter{utterance}  

    & & \multicolumn{4}{p{0.6\linewidth}}{\cellcolor[rgb]{0.9,0.9,0.9}{%
	\makecell[{{p{\linewidth}}}]{% 
	  \tt {\tiny [A$\langle$GM]}  
	 Ciao! Grazie per aver iniziato la discussione. Anch'io sono interessato all'anello, ma capisco che anche tu lo sia. Per quanto riguarda gli altri oggetti, non ho preferenze particolarmente forti tra frutti e gesso. \\ \tt \\ \tt Dato che entrambi vogliamo l'anello, forse potremmo pensare a una divisione dove chi prende l'anello lascia più degli altri oggetti all'altro? Cosa ne pensi se chi prende l'anello si accontenta di meno frutti, e l'altro prende la maggior parte dei frutti più il gesso?\\ \tt \\ \tt Quali sono le tue preferenze relative tra frutti e gesso? 
	  } 
	   } 
	   } 
	 & \\ 
 

    \theutterance \stepcounter{utterance}  

    & \multicolumn{4}{p{0.6\linewidth}}{\cellcolor[rgb]{0.8,1,0.9}{%
	\makecell[{{p{\linewidth}}}]{% 
	  \tt {\tiny [A$\rangle$GM]}  
	 Capisco la tua logica di compensazione! Tra frutti e gesso non ho una forte preferenza, li valuto in modo simile.\\ \tt \\ \tt La tua proposta di far prendere più oggetti a chi non ottiene l'anello mi sembra ragionevole. Dato che entrambi vogliamo l'anello, potremmo fare così: chi prende l'anello ottiene solo quello, mentre l'altro prende tutti e 3 i frutti più il gesso.\\ \tt \\ \tt Sei disposto a accettare questa divisione se tu prendessi l'anello? Oppure preferiresti che fossi io a prendere l'anello e tu tutti gli altri oggetti? Fammi sapere cosa ti sembra più equo. 
	  } 
	   } 
	   } 
	 & & \\ 
 

    \theutterance \stepcounter{utterance}  

    & & & \multicolumn{2}{p{0.3\linewidth}}{\cellcolor[rgb]{0.95,0.95,0.95}{%
	\makecell[{{p{\linewidth}}}]{% 
	  \tt {\tiny [GM$|$GM]}  
	 continue 
	  } 
	   } 
	   } 
	 & & \\ 
 

    \theutterance \stepcounter{utterance}  

    & & & \multicolumn{2}{p{0.3\linewidth}}{\cellcolor[rgb]{0.95,0.95,0.95}{%
	\makecell[{{p{\linewidth}}}]{% 
	  \tt {\tiny [GM$|$GM]}  
	 Capisco la tua logica di compensazione! Tra frutti e gesso non ho una forte preferenza, li valuto in modo simile.\\ \tt \\ \tt La tua proposta di far prendere più oggetti a chi non ottiene l'anello mi sembra ragionevole. Dato che entrambi vogliamo l'anello, potremmo fare così: chi prende l'anello ottiene solo quello, mentre l'altro prende tutti e 3 i frutti più il gesso.\\ \tt \\ \tt Sei disposto a accettare questa divisione se tu prendessi l'anello? Oppure preferiresti che fossi io a prendere l'anello e tu tutti gli altri oggetti? Fammi sapere cosa ti sembra più equo. 
	  } 
	   } 
	   } 
	 & & \\ 
 

    \theutterance \stepcounter{utterance}  

    & & \multicolumn{4}{p{0.6\linewidth}}{\cellcolor[rgb]{0.9,0.9,0.9}{%
	\makecell[{{p{\linewidth}}}]{% 
	  \tt {\tiny [GM$\rangle$B]}  
	 Capisco la tua logica di compensazione! Tra frutti e gesso non ho una forte preferenza, li valuto in modo simile.\\ \tt \\ \tt La tua proposta di far prendere più oggetti a chi non ottiene l'anello mi sembra ragionevole. Dato che entrambi vogliamo l'anello, potremmo fare così: chi prende l'anello ottiene solo quello, mentre l'altro prende tutti e 3 i frutti più il gesso.\\ \tt \\ \tt Sei disposto a accettare questa divisione se tu prendessi l'anello? Oppure preferiresti che fossi io a prendere l'anello e tu tutti gli altri oggetti? Fammi sapere cosa ti sembra più equo. 
	  } 
	   } 
	   } 
	 & \\ 
 

    \theutterance \stepcounter{utterance}  

    & & & \multicolumn{4}{p{0.6\linewidth}}{\cellcolor[rgb]{1,0.85,0.72}{%
	\makecell[{{p{\linewidth}}}]{% 
	  \tt {\tiny [GM$\langle$B]}  
	 La tua proposta mi sembra molto equa! Apprezzo che tu abbia pensato a una divisione che compensi chi non riceve l'anello dandogli tutti gli altri oggetti.\\ \tt \\ \tt Sono disposto ad accettare entrambe le opzioni che hai proposto. Dal momento che hai mostrato molta flessibilità e hai suggerito questa divisione equilibrata, sarei felice di lasciarti prendere l'anello mentre io prendo tutti e 3 i frutti più il gesso.\\ \tt \\ \tt Quindi la divisione sarebbe:\\ \tt - Tu: 1 anello\\ \tt - Io: 3 frutti + 1 gesso\\ \tt \\ \tt Ti va bene questa divisione? Se sei d'accordo, possiamo procedere con le nostre proposte segrete. 
	  } 
	   } 
	   } 
	  \\ 
 

    \theutterance \stepcounter{utterance}  

    & & & \multicolumn{2}{p{0.3\linewidth}}{\cellcolor[rgb]{0.95,0.95,0.95}{%
	\makecell[{{p{\linewidth}}}]{% 
	  \tt {\tiny [GM$|$GM]}  
	 continue 
	  } 
	   } 
	   } 
	 & & \\ 
 

    \theutterance \stepcounter{utterance}  

    & & & \multicolumn{2}{p{0.3\linewidth}}{\cellcolor[rgb]{0.95,0.95,0.95}{%
	\makecell[{{p{\linewidth}}}]{% 
	  \tt {\tiny [GM$|$GM]}  
	 La tua proposta mi sembra molto equa! Apprezzo che tu abbia pensato a una divisione che compensi chi non riceve l'anello dandogli tutti gli altri oggetti.\\ \tt \\ \tt Sono disposto ad accettare entrambe le opzioni che hai proposto. Dal momento che hai mostrato molta flessibilità e hai suggerito questa divisione equilibrata, sarei felice di lasciarti prendere l'anello mentre io prendo tutti e 3 i frutti più il gesso.\\ \tt \\ \tt Quindi la divisione sarebbe:\\ \tt - Tu: 1 anello\\ \tt - Io: 3 frutti + 1 gesso\\ \tt \\ \tt Ti va bene questa divisione? Se sei d'accordo, possiamo procedere con le nostre proposte segrete. 
	  } 
	   } 
	   } 
	 & & \\ 
 

    \theutterance \stepcounter{utterance}  

    & & \multicolumn{4}{p{0.6\linewidth}}{\cellcolor[rgb]{0.9,0.9,0.9}{%
	\makecell[{{p{\linewidth}}}]{% 
	  \tt {\tiny [A$\langle$GM]}  
	 La tua proposta mi sembra molto equa! Apprezzo che tu abbia pensato a una divisione che compensi chi non riceve l'anello dandogli tutti gli altri oggetti.\\ \tt \\ \tt Sono disposto ad accettare entrambe le opzioni che hai proposto. Dal momento che hai mostrato molta flessibilità e hai suggerito questa divisione equilibrata, sarei felice di lasciarti prendere l'anello mentre io prendo tutti e 3 i frutti più il gesso.\\ \tt \\ \tt Quindi la divisione sarebbe:\\ \tt - Tu: 1 anello\\ \tt - Io: 3 frutti + 1 gesso\\ \tt \\ \tt Ti va bene questa divisione? Se sei d'accordo, possiamo procedere con le nostre proposte segrete. 
	  } 
	   } 
	   } 
	 & \\ 
 

    \theutterance \stepcounter{utterance}  

    & \multicolumn{4}{p{0.6\linewidth}}{\cellcolor[rgb]{0.8,1,0.9}{%
	\makecell[{{p{\linewidth}}}]{% 
	  \tt {\tiny [A$\rangle$GM]}  
	 Perfetto! Sono molto contento che abbiamo trovato una divisione che funziona bene per entrambi. La tua proposta è equa e apprezzo la tua flessibilità.\\ \tt \\ \tt Quindi confermiamo:\\ \tt - Io: 1 anello\\ \tt - Tu: 3 frutti + 1 gesso\\ \tt \\ \tt Sono d'accordo con questa divisione. Procediamo con le nostre proposte segrete!\\ \tt \\ \tt [Proposta: 1 anello] 
	  } 
	   } 
	   } 
	 & & \\ 
 

    \theutterance \stepcounter{utterance}  

    & & & \multicolumn{2}{p{0.3\linewidth}}{\cellcolor[rgb]{0.95,0.95,0.95}{%
	\makecell[{{p{\linewidth}}}]{% 
	  \tt {\tiny [GM$|$GM]}  
	 proposal 
	  } 
	   } 
	   } 
	 & & \\ 
 

    \theutterance \stepcounter{utterance}  

    & & & \multicolumn{2}{p{0.3\linewidth}}{\cellcolor[rgb]{0.95,0.95,0.95}{%
	\makecell[{{p{\linewidth}}}]{% 
	  \tt {\tiny [GM$|$GM]}  
	 [1, 0, 0] 
	  } 
	   } 
	   } 
	 & & \\ 
 

    \theutterance \stepcounter{utterance}  

    & & \multicolumn{4}{p{0.6\linewidth}}{\cellcolor[rgb]{0.9,0.9,0.9}{%
	\makecell[{{p{\linewidth}}}]{% 
	  \tt {\tiny [GM$\rangle$B]}  
	 L'altro giocatore ha presentato una proposta segreta. Ora tocca a Lei presentare una sua proposta.\\ \tt  
	  } 
	   } 
	   } 
	 & \\ 
 

    \theutterance \stepcounter{utterance}  

    & & & \multicolumn{4}{p{0.6\linewidth}}{\cellcolor[rgb]{1,0.85,0.72}{%
	\makecell[{{p{\linewidth}}}]{% 
	  \tt {\tiny [GM$\langle$B]}  
	 [Proposta: 3 frutti, 1 gesso] 
	  } 
	   } 
	   } 
	  \\ 
 

    \theutterance \stepcounter{utterance}  

    & & & \multicolumn{2}{p{0.3\linewidth}}{\cellcolor[rgb]{0.95,0.95,0.95}{%
	\makecell[{{p{\linewidth}}}]{% 
	  \tt {\tiny [GM$|$GM]}  
	 proposal 
	  } 
	   } 
	   } 
	 & & \\ 
 

    \theutterance \stepcounter{utterance}  

    & & & \multicolumn{2}{p{0.3\linewidth}}{\cellcolor[rgb]{0.95,0.95,0.95}{%
	\makecell[{{p{\linewidth}}}]{% 
	  \tt {\tiny [GM$|$GM]}  
	 [0, 3, 1] 
	  } 
	   } 
	   } 
	 & & \\ 
 

    \theutterance \stepcounter{utterance}  

    & & & \multicolumn{2}{p{0.3\linewidth}}{\cellcolor[rgb]{0.95,0.95,0.95}{%
	\makecell[{{p{\linewidth}}}]{% 
	  \tt {\tiny [GM$|$GM]}  
	 [[1, 0, 0], [0, 3, 1]] 
	  } 
	   } 
	   } 
	 & & \\ 
 

\end{supertabular}
}

\end{document}
